% =====================================================================
%  preamble.tex —— 书籍信息与正文用到的少量宏
%
%  排版定义（定理类环境、证明、列表、编号）在 qyxf-book.cls 中；
%  OCR 导出的 Unicode 字符映射在 preamble/_unicode-map.tex 中。
%  本文件只放这两样之外的东西。
% =====================================================================

% 不要载入 amssymb：文类已用 unicode-math 载入数学字体，两者对
% \square、\blacksquare、\eth、\smallsetminus、\digamma、\backepsilon
% 等符号重复定义，会直接报 "Command ... already defined"。
% unicode-math 缺的几个同义符号在 qyxf-book.cls 里已用 \providecommand 补齐；
% \xrightarrow / \xleftarrow 由 LaTeX2e 内核提供，无需额外宏包。
\usepackage{stmaryrd}          % \rightleftarrows：原文用 ⇄ 标注证明方向
\usepackage{newunicodechar}
\usepackage{array,makecell}
\usepackage{indentfirst}
\usepackage{fvextra}           % damagedformula：无法排版的公式原样展示
\usepackage{needspace}         % A2：小标题/带框条目起排前的最小空间保护（避免页底孤行）
% 交换图。OCR 把原书的交换图拆成了散落的 array + \downarrow 文字，
% 版面完全走形；这里用 tikz-cd 重画（见项目文档《交换图与表格规范》）。
\usepackage{tikz-cd}
\usetikzlibrary{arrows.meta}
\tikzcdset{arrow style=tikz, diagrams={>={Straight Barb[scale=0.8]}}}

% 数学粗体字母（\symbfup 等）取自 XITS Math Bold；不声明的话 unicode-math 会把
% 粗体数学字母回退到正文字体 XITS Bold，希腊字母（U+1D6E9 一带）在那里是缺字。
\setmathfont{XITS Math Bold}[version=bold]

% ------------------------------------------------------------------ 版式
% 首行缩进两字、段间不空行、行距按 ctex 中文默认；应急伸长留足，
% 以免 OCR 造成的长串不可断行文字冲出右边界。
\setlength{\emergencystretch}{3em}
\sloppy
\hyphenpenalty=10000
\exhyphenpenalty=10000

% 表格：OCR 真值表按 booktabs 排出，字号略微收小以便与正文区分
\renewcommand{\arraystretch}{1.15}

% ------------------------------------------------------------------ 宏
% 缺图占位（图片缓存缺失时使用）
\newcommand{\figplaceholder}[2]{%
  \fbox{\begin{minipage}[c][3.2cm][c]{#1\textwidth}\centering\small\itshape [ #2 ]\end{minipage}}}
% 无法排版的公式：原样展示（OCR 已损坏的公式，按项目铁律不改写源文）
\DefineVerbatimEnvironment{damagedformula}{Verbatim}{breaklines=true,breakanywhere=true,fontsize=\scriptsize,frame=single,framesep=4pt}

% 图：原文的插图居中放置，不编号（原文也未给插图编号）
\newcommand{\bookfigure}[2][0.5]{%
  \begin{center}\includegraphics[width=#1\textwidth]{#2}\end{center}}

% ------------------------------------------------------------ 书籍信息
\title{分析学 II}
\subtitle{Analysis II}
\author{Herbert Amann \quad Joachim Escher}
\date{2026 年 10 月}
\org{\textit{中文译本}\\[1ex]\textsc{Analysis II, Amann \& Escher}}
\footorg{分析学 II · 中文译本}
\license{%
本译本为原著的中文翻译，仅供个人学习与研究使用，不得用于任何商业用途。\\[1ex]
原著 \textit{Analysis II}（Herbert Amann, Joachim Escher）版权归作者与
Birkhäuser Verlag 所有，英译由 Silvio Levy 与 Matthew Cargo 完成（Birkhäuser Verlag，2008）。\\[1ex]
翻译、排版：Jea神纳}
